% AA-29 | Iteracja XV; szczegółowy audyt historyczny III, 20 XI 2026.
\chapter[Wielomiany nad ciałami skończonymi]{Struktura ciał skończonych i wielomiany nierozkładalne}
\label{aa:ch:finite-fields-structure}

\section{Frobenius i rozkład wielomianów}

Znajomość liczby elementów ciała skończonego daje narzędzie do
znajdowania jego elementów algebraicznych. Jeżeli ciało ma $q$
elementów, każdy jego element jest pierwiastkiem $X^q-X$.
Uogólnienie do $X^{q^n}-X$ pozwoli scharakteryzować wszystkie
moniczne wielomiany nierozkładalne, których stopnie dzielą $n$.
Należy przy tym odróżnić dwa stwierdzenia: wielomian
$X^{q^n}-X$ jest rozdzielczy, lecz dowolny wielomian
stopnia $n$ nad ciałem $q$-elementowym nie musi być rozdzielczy
ani nierozkładalny.

\begin{twierdzenie}[Pierwiastki wielomianu Frobeniusa]
\label{aa:tw:finite-frobenius-roots}
Niech $p$ będzie liczbą pierwszą, $m,n\ge1$ liczbami
całkowitymi, a $q=p^m$. W ustalonym domknięciu
algebraicznym $\overline{\mathbb F}_p$ połóżmy
\[
E_q=\{x\in\overline{\mathbb F}_p:x^q=x\}.
\]
Wówczas $E_q$ jest ciałem o dokładnie $q$ elementach.
Ponadto w pierścieniu $E_q[X]$ zachodzi równość
\[
X^{q^n}-X=\prod_{\substack{f\in E_q[X]\ \text{$f$ moniczny, nierozkładalny}\\ \deg f\mid n}}f(X),
\]
gdzie każdy czynnik występuje dokładnie jeden raz.
\end{twierdzenie}
\begin{proof}
\emph{Krok 1: ciało $E_q$.} W charakterystyce $p$
iteracja tożsamości $(x+y)^p=x^p+y^p$ daje
$(x+y)^q=x^q+y^q$ dla $q=p^m$.
Wobec tego, jeśli $x^q=x$ oraz $y^q=y$, to
$(x+y)^q=x+y$ i $(xy)^q=xy$.
Ponadto $(-x)^q=-x$ (w charakterystyce $2$ mamy
$-x=x$), a dla $x\ne0$ również
$(x^{-1})^q=(x^q)^{-1}=x^{-1}$.
Elementy $0$ i $1$ należą do $E_q$, więc wszystkie
aksjomaty podciała są spełnione.
Pochodna $X^q-X$ jest równa $-1$, ponieważ
$q=0$ w $\mathbb F_p$; wielomian jest więc rozdzielczy.
W domknięciu algebraicznym ma dokładnie $q$ różnych
pierwiastków, czyli $|E_q|=q$.

\emph{Krok 2: ciało o $q^d$ elementach.}
Niech $f\in E_q[X]$ będzie moniczny i nierozkładalny
stopnia $d\ge1$, a $\alpha\in\overline{\mathbb F}_p$
jego pierwiastkiem. Iloraz $E_q[X]/(f)$ jest ciałem
oraz ma bazę $1,\bar X,\ldots,\bar X^{d-1}$ nad
$E_q$. Wobec tego ma dokładnie $q^d$ elementów.
Ewaluacja $\bar X\mapsto\alpha$ identyfikuje ten
iloraz z $E_q(\alpha)$, zatem
$|E_q(\alpha)|=q^d$.
Jego grupa multiplikatywna ma $q^d-1$ elementów,
więc dla $\alpha\ne0$ twierdzenie Lagrange'a daje
$\alpha^{q^d-1}=1$, a stąd
$\alpha^{q^d}=\alpha$. Dla $\alpha=0$ ta sama
równość jest natychmiastowa.

\emph{Krok 3: minimalny okres Frobeniusa.}
Dla każdej liczby całkowitej $r\ge1$ oznaczmy
$\operatorname{Fr}_q^r(x)=x^{q^r}$, gdzie
$\operatorname{Fr}_q:x\mapsto x^q$ jest iteracją
potęgowania; na każdym ciele $E_{q^s}$ jest automorfizmem,
ponieważ zachowuje $E_{q^s}$ (dla $x^{q^s}=x$ mamy
$(x^q)^{q^s}=x^q$) i jest iniektywny:
z $x^q=y^q$ wynika $(x-y)^q=0$, a więc $x=y$.
Ponieważ $E_{q^s}$ jest skończone, potęgowanie przez
$q$ jest również surjektywne.
Niech $r$ będzie najmniejszą dodatnią liczbą taką, że
$\alpha^{q^r}=\alpha$. Istnieje ona i $r\le d$, na
mocy poprzedniego kroku.
Ponieważ $\alpha\in E_{q^r}$ oraz $E_q\subset E_{q^r}$
(jeśli $x^q=x$, to kolejne potęgowania przez $q$
także pozostawiają $x$ nieruchomy),
$E_q(\alpha)$ jest podciałem $E_{q^r}$.
Porównanie liczności podciała i ciała daje
$q^d\le q^r$, czyli $d\le r$. Łącznie z
$r\le d$ wynika $r=d$.

Teraz wykażemy, że
$\alpha^{q^n}=\alpha$ wtedy i tylko wtedy, gdy $d\mid n$.
Kierunek $d\mid n$ wynika z wielokrotnego zastosowania
$\operatorname{Fr}_q^d(\alpha)=\alpha$.
W kierunku odwrotnym podzielmy $n$ z resztą przez $d$:
$n=ad+s$ dla $a\ge0$ i $0\le s<d$.
Z uwagi na $\operatorname{Fr}_q^d(\alpha)=\alpha$
mamy $\operatorname{Fr}_q^n(\alpha)=
\operatorname{Fr}_q^s(\alpha)$.
Jeżeli lewa strona jest równa $\alpha$, to minimalność
$d$ wymusza $s=0$, czyli $d\mid n$.

\emph{Krok 4: identyfikacja wszystkich czynników.}
Wielomian moniczny nierozkładalny $f$ dzieli
$X^{q^n}-X$ dokładnie wtedy, gdy jego pierwiastek
$\alpha$ spełnia $\alpha^{q^n}=\alpha$, a więc
wtedy i tylko wtedy, gdy $\deg f=d$ dzieli $n$.
Pochodna $X^{q^n}-X$ jest równa $-1$, toteż
w jego rozkładzie nie ma powtórzonych czynników.
Ponieważ każdy niestały wielomian nad ciałem rozkłada się
na czynniki nierozkładalne, otrzymujemy wskazaną
faktoryzację.
\end{proof}

\begin{wniosek}[Tożsamość stopni]\label{aa:cor:finite-degree-count}
Niech $q=p^m$, gdzie $p$ jest liczbą pierwszą i $m\ge1$.
Dla każdej liczby całkowitej $d\ge1$ oznaczmy przez
$N_q(d)$ liczbę monicznych wielomianów nierozkładalnych
stopnia $d$ w $E_q[X]$. Dla każdego $n\ge1$ zachodzi
\[
q^n=\sum_{d\mid n}d\,N_q(d).
\]
\end{wniosek}
\begin{proof}
Po lewej stronie faktoryzacji z
Twierdzenia~\ref{aa:tw:finite-frobenius-roots} znajduje się
wielomian stopnia $q^n$. Dla każdego dzielnika $d\mid n$
po prawej stronie występuje dokładnie $N_q(d)$ czynników
stopnia $d$. Stopień iloczynu niezerowych wielomianów
nad ciałem jest sumą stopni czynników, skąd wzór.
\end{proof}

\section{Odwrócenie Möbiusa i liczba wielomianów}

Potrzebujemy teraz odwrócić sumę po dzielnikach $n$.
Umożliwia to funkcja Möbiusa: jest ona dobrana tak, aby
jej suma po wszystkich dzielnikach była zerowa poza
przypadkiem $n=1$.

\begin{definicja}\label{aa:def:mobius}
Dla każdej liczby całkowitej $n\ge1$ definiujemy
funkcję Möbiusa $\mu(n)$ następująco:
$\mu(1)=1$; jeśli dla pewnej liczby pierwszej $p$
zachodzi $p^2\mid n$, to $\mu(n)=0$;
w pozostałym przypadku $n>1$ jest iloczynem $r\ge1$
\emph{różnych} liczb pierwszych i $\mu(n)=(-1)^r$.
\end{definicja}
\begin{lemat}\label{aa:lem:mobius-sum}
Dla każdej liczby całkowitej $n\ge1$ mamy
\[
\sum_{d\mid n}\mu(d)=
\begin{cases}1,&n=1,\\0,&n>1.\end{cases}
\]
\end{lemat}
\begin{proof}
Dla $n=1$ jedynym dzielnikiem dodatnim jest $1$.
Jeżeli $n>1$, niech $p_1,\ldots,p_r$ będą jego
różnymi dzielnikami pierwszymi, gdzie $r\ge1$.
Składnik $\mu(d)$ jest niezerowy dokładnie dla tych
dzielników $d$, które są iloczynami pewnego podzbioru
zbioru $\{p_1,\ldots,p_r\}$, w tym podzbioru pustego.
Podzbiór $j$-elementowy wnosi $(-1)^j$, a takich
podzbiorów jest $\binom rj$. Otrzymujemy
$\sum_{d\mid n}\mu(d)
=\sum_{j=0}^r\binom rj(-1)^j=(1-1)^r=0$.
\end{proof}

\begin{twierdzenie}[Wzór na liczbę wielomianów nierozkładalnych]
\label{aa:tw:finite-irreducible-count}
Niech $p$ będzie liczbą pierwszą, $m,n\ge1$ liczbami
całkowitymi, $q=p^m$, a $E_q$ ciałem o $q$ elementach.
Niech $N_q(n)$ oznacza liczbę monicznych nierozkładalnych
wielomianów stopnia $n$ w $E_q[X]$. Wówczas
\[
N_q(n)=\frac1n\sum_{d\mid n}\mu(d)q^{n/d}.
\]
Ponadto $N_q(n)>0$: nad $E_q$ istnieje moniczny
wielomian nierozkładalny każdego dodatniego stopnia $n$.
\end{twierdzenie}
\begin{proof}
\emph{Krok 1: odwrócenie sumy po dzielnikach.}
Z Wniosku~\ref{aa:cor:finite-degree-count} dla dowolnego
$k\ge1$ mamy $q^k=\sum_{e\mid k}eN_q(e)$.
Wstawiając tę tożsamość do sumy z tezy, uzyskujemy
\[
\begin{aligned}
\sum_{d\mid n}\mu(d)q^{n/d}
&=\sum_{d\mid n}\mu(d)\sum_{e\mid n/d}eN_q(e)\\
&=\sum_{e\mid n}eN_q(e)\sum_{d\mid n/e}\mu(d)\\
&=nN_q(n).
\end{aligned}
\]
Druga równość jest dopuszczalna, ponieważ sumy są
skończone, a para $(d,e)$ spełnia $de\mid n$ dokładnie
wtedy, gdy $e\mid n$ i $d\mid n/e$. Ostatnia równość
wynika z Lematu~\ref{aa:lem:mobius-sum}: suma wewnętrzna
jest równa $1$ tylko dla $e=n$.
Dzieląc obie strony przez dodatnią liczbę $n$,
otrzymujemy żądany wzór.

\emph{Krok 2: istnienie nierozkładalnego wielomianu.}
Dla $n=1$ wystarczy moniczny wielomian liniowy $X$.
Jest on nierozkładalny nad każdym ciałem.
Niech $n\ge2$. Ciało $E_{q^n}$ ma $q^n$ elementów
i zawiera $E_q$. Rozważmy zbiór tych elementów
$\alpha\in E_{q^n}$, których stopień
$[E_q(\alpha):E_q]$ jest właściwym dzielnikiem $n$.
Jest to rzeczywiście dzielnik $n$ na mocy twierdzenia
o wieży rozszerzeń, bo $E_q\subseteq E_q(\alpha)
\subseteq E_{q^n}$ i $[E_{q^n}:E_q]=n$: ostatnia
równość wynika z $|E_{q^n}|=q^n$ i $|E_q|=q$.
Jeśli ten stopień jest równy $d$, to
$\alpha\in E_{q^d}$, ponieważ każde ciało o $q^d$
elementach spełnia równanie $x^{q^d}=x$.
Zatem liczba wszystkich takich elementów jest
co najwyżej
\[
\sum_{\substack{d\mid n\\d<n}}q^d
\le \sum_{j=1}^{\lfloor n/2\rfloor}q^j
<q^n.
\]
Pierwsza nierówność wynika z faktu, że każdy właściwy
dodatni dzielnik $n$ jest co najwyżej $n/2$.
Druga jest ścisła: dla $k\ge1$ i $q\ge2$
\(\sum_{j=1}^{k}q^j
=q(q^k-1)/(q-1)<q^{k+1}\), a
$\lfloor n/2\rfloor+1\le n$.
Istnieje wobec tego $\alpha\in E_{q^n}$, którego
stopień nad $E_q$ nie jest właściwym dzielnikiem $n$.
Jest to z konieczności stopień $n$. Wielomian minimalny
$\alpha$ nad $E_q$ jest moniczny, nierozkładalny i ma
stopień $n$. Stąd $N_q(n)\ge1$.
\end{proof}

\begin{przyklad}[Stopień dwa nad $\mathbb F_2$]
\label{aa:prz:finite-quadratic}
Dla $q=2$ i $n=2$ wzór daje
$N_2(2)=(2^2-2)/2=1$. Wielomian
$f=X^2+X+1$ nie ma pierwiastka w $\mathbb F_2$:
$f(0)=1$, $f(1)=1+1+1=1$.
Wielomian stopnia $2$ bez pierwiastka jest
nierozkładalny, ponieważ ewentualny rozkład na
czynniki niestałe musiałby być iloczynem dwóch
wielomianów liniowych. Istnieje więc dokładnie jeden
moniczny nierozkładalny wielomian kwadratowy.
Iloraz $\mathbb F_2[X]/(f)$ jest ciałem i jako
$\mathbb F_2$-przestrzeń liniowa ma bazę $1,x$,
gdzie $x$ oznacza klasę $X$. Jego elementami są
$0,1,x,x+1$. Równość $f(x)=0$ daje
$x^2=x+1$, a więc $x^3=x(x+1)=x^2+x=1$.
Dla kontroli rozkładu Frobeniusa mamy
\[
X^4-X=X(X+1)(X^2+X+1)
\quad\text{w }\mathbb F_2[X].
\]
Iloczyn po prawej stronie ma stopień $4$ i można go
sprawdzić przez wymnożenie; odpowiada liniowym
czynnikom stopnia $1$ i jednemu czynnikowi stopnia $2$.
\end{przyklad}

\begin{przyklad}[Kontrprzykład: nie każdy wielomian stopnia $n$ jest nierozkładalny]
\label{aa:prz:finite-reducible}
W $\mathbb F_2[X]$ zachodzi
$X^4+X^2+1=(X^2+X+1)^2$. Jest to wielomian
rozkładalny, którego pochodna wynosi zero.
W pierścieniu
$R=\mathbb F_2[X]/(X^4+X^2+1)$
klasa $v$ wielomianu $X^2+X+1$ spełnia $v^2=0$.
Jest niezerowa: gdyby należała do ideału generowanego
przez wielomian stopnia $4$, byłaby równa zeru albo
miałaby stopień co najmniej $4$, podczas gdy ma stopień
$2$. Zatem $R$ zawiera niezerowy element nilpotentny.
Żadne ciało nie ma takiego elementu, bo z $v^2=0$
i $v\ne0$ można by pomnożyć przez $v^{-1}$ i uzyskać
$v=0$. Chociaż $R$ ma $2^4=16$ elementów jako
czterowymiarowa przestrzeń nad $\mathbb F_2$, nie jest
ciałem szesnastoelementowym.
\end{przyklad}

\begin{przyklad}[Wielomiany trzeciego stopnia nad $\mathbb F_2$]
\label{aa:prz:finite-cubics}
Dla $n=3$ wzór Möbiusa daje
$N_2(3)=(2^3-2)/3=2$.
Sprawdzimy to niezależnie. Wielomian stopnia $3$
nad ciałem jest nierozkładalny dokładnie wtedy,
gdy nie ma pierwiastka w tym ciele: jeśli jest
rozkładalny, stopnie jego niestałych czynników sumują
się do $3$, więc jeden z nich ma stopień $1$.
Każdy moniczny wielomian sześcienny w $\mathbb F_2[X]$
ma postać $X^3+aX^2+bX+c$ dla $a,b,c\in\{0,1\}$.
Aby nie miał pierwiastka $0$, musi być $c=1$.
Wtedy jego wartość dla $X=1$ to
$1+a+b+1=a+b$, więc brak pierwiastka $1$ wymaga
$a+b=1$. Są dokładnie dwie pary współczynników
$(a,b)$ spełniające tę równość: $(0,1)$ i $(1,0)$.
Otrzymujemy wielomiany
$X^3+X+1$ i $X^3+X^2+1$, oba nierozkładalne.
Wszystkie pozostałe wielomiany stopnia $3$ mają
pierwiastek w $\mathbb F_2$, a zatem są rozkładalne.
Istotnie, zgodnie z Twierdzeniem~\ref{aa:tw:finite-frobenius-roots},
\[
X^8-X=X(X+1)(X^3+X+1)(X^3+X^2+1)
\quad\text{w }\mathbb F_2[X].
\]
\end{przyklad}

\begin{przyklad}[Trzy wielomiany kwadratowe nad $\mathbb F_3$]
\label{aa:prz:finite-quadratics-three}
Wzór dla $q=3$ i $n=2$ daje
$N_3(2)=(9-3)/2=3$.
Moniczny wielomian kwadratowy ma postać
$X^2+aX+b$ z $a,b\in\mathbb F_3$.
Jest nierozkładalny dokładnie wtedy, gdy jego
wyróżnik $\Delta=a^2-4b=a^2-b$ nie jest kwadratem
w $\mathbb F_3$: istotnie, ponieważ $2$ jest
odwracalne, równanie $X^2+aX+b=0$ jest równoważne
$(2X+a)^2=\Delta$. Ponieważ kwadratami są tylko $0$ i $1$,
warunek ten jest równoważny $\Delta=2$.
Jeśli $a=0$, dostajemy $b=1$ i $X^2+1$.
Jeśli $a=1$ lub $a=2$, to $a^2=1$, więc $b=2$,
co daje $X^2+X+2$ oraz $X^2+2X+2$.
W ten sposób niezależnie znaleźliśmy wszystkie trzy
wielomiany nierozkładalne stopnia $2$ nad $\mathbb F_3$.
\end{przyklad}

\begin{uwaga}\label{aa:uw:finite-field-previous}
Istnienie i jednoznaczność ciał skończonych były już
omawiane w Rozdziale~\ref{aa:ch:embeddings-normal-separable}.
Tutaj rozwinęliśmy rozkład wielomianów Frobeniusa,
ścisłe wyliczanie liczby wielomianów nierozkładalnych
oraz konstrukcję elementów stopnia dokładnie $n$.
\end{uwaga}
